\section{Relacionamentos e regras de negócio}

Descrevam cada relação usando um verbo e anotem o que precisa ser confirmado no processo real.

\begin{table}[!ht]
\centering
\small
\setlength{\tabcolsep}{4pt}
\renewcommand{\arraystretch}{1.75}
\begin{tabularx}{\linewidth}{|>{\raggedright\arraybackslash}X
                              |>{\centering\arraybackslash}X
                              |>{\raggedright\arraybackslash}X
                              |>{\raggedright\arraybackslash}X|}
\hline
\textbf{Entidade A} &
\textbf{Relação (verbo)} &
\textbf{Entidade B} &
\textbf{Regra / dúvida a validar} \\
\hline
\textbf{Estudante} & realiza & \textbf{Solicitação} &
Quantas solicitações um estudante pode realizar? \\
\hline
 & & & \\ \hline
 & & & \\ \hline
 & & & \\ \hline
 & & & \\ \hline
 & & & \\ \hline
 & & & \\ \hline
 & & & \\ \hline
 & & & \\ \hline
 & & & \\ \hline
 & & & \\ \hline
 & & & \\ \hline
\end{tabularx}
\end{table}

\begin{example}{Exemplo de continuidade}{continuidademer}
O BPMN registra a análise da solicitação, um requisito permite registrar a pendência
e os dados identificados são \emph{solicitação} e \emph{pendência}. Assim, no MER,
\textbf{SOLICITAÇÃO possui PENDÊNCIA}. A quantidade de pendências depende da regra real do processo.
\end{example}

\begin{warnbox}[Observação]
Se faltar linha, continue registrando os relacionamentos na parte de trás da folha ou em uma folha adicional.
\end{warnbox}
